% 原分章PDF83页附录7.A标题—88页附录7.B标题之前，印刷页659—664。
\section*{附录7.A\quad LiNGAM的对数似然比近似}
\addcontentsline{toc}{section}{附录7.A LiNGAM的对数似然比近似}
\markright{附录7.A LiNGAM的对数似然比近似}

近似式的推导采用Hyvärinen and Smith（2013）的方法。考虑式（7.3）中的因果模型 $X\to Y$：
\begin{equation}
\begin{aligned}
X&=\alpha u+\varepsilon_x,\\
Y&=\rho X+\beta u+\varepsilon_y.
\end{aligned}\tag{7.A.1}
\end{equation}
首先，对变量 $\varepsilon_Y$ 进行标准化。由式（7.A.1），有
\begin{equation}
\varepsilon_Y=Y-\rho X-\beta u,\tag{7.A.2}
\end{equation}
从而
\begin{equation}
\operatorname{Var}(\varepsilon_Y)=\operatorname{Var}(Y)
 +\operatorname{Var}(\rho X)-2\operatorname{cov}(Y,\rho X).\tag{7.A.3}
\end{equation}
假设变量 $X,Y$ 已经标准化，则有
\begin{equation}
\begin{aligned}
\operatorname{Var}(X)&=1,\\
\operatorname{Var}(Y)&=1,\\
\operatorname{Cov}(X,Y)&=\rho\operatorname{Var}(X)=\rho.
\end{aligned}\tag{7.A.4}
\end{equation}
将式（7.A.4）代入式（7.A.3），得到
\begin{equation}
\begin{aligned}
\operatorname{Var}(\varepsilon_Y)
 &=\operatorname{Var}(Y)+\rho^2\operatorname{Var}(X)
   -2\rho\operatorname{cov}(Y,X)\\
 &=1+\rho^2-2\rho^2=1-\rho^2.
\end{aligned}\tag{7.A.5}
\end{equation}
定义标准化后的 $\varepsilon_Y$ 为
\begin{equation}
d=\frac{\varepsilon_Y}{\sqrt{1-\rho^2}}.\tag{7.A.6}
\end{equation}

令 $P_x(u)$ 为随机变量 $X$ 的密度函数，$P_d(u)$ 为变量 $d$ 的密度函数。由式（7.A.1），$\varepsilon_X$ 的密度函数等于变量 $X$ 的密度函数 $P_x(u)$。由式（7.A.6），显然有
\begin{equation}
P(\varepsilon_Y)=\frac{P_d(u)}{\sqrt{1-\rho^2}}.\tag{7.A.7}
\end{equation}
定义 $G_x(u)=\log P_x(u)$、$G_d(u)=\log P_d(u)$。随机变量 $X,Y$ 的联合密度函数为
\begin{equation}
\begin{aligned}
P(x,y)&=P(\varepsilon_X,\varepsilon_Y)\\
 &=P(\varepsilon_X)P(\varepsilon_Y).
\end{aligned}\tag{7.A.8}
\end{equation}
将式（7.A.7）代入式（7.A.8），得到
\begin{equation}
P(x,y)=P_x(x)\frac{P_d\left(\dfrac{y-\rho x-\beta u}{\sqrt{1-\rho^2}}\right)}
 {\sqrt{1-\rho^2}}.\tag{7.A.9}
\end{equation}
对式（7.A.9）两边取对数，得到式（7.A.1）中 $X\to Y$ 的LiNGAM的对数似然：
\begin{equation}
l(X\to Y)=\sum_{i=1}^n\left[G_x(x_i)+G_d\left(
 \frac{y_i-\rho x_i-\beta u_i}{\sqrt{1-\rho^2}}\right)\right]
 -\frac n2\log(1-\rho^2).\tag{7.A.10}
\end{equation}
回顾式（7.4）中的第二个潜在因果模型 $Y\to X$：
\begin{equation}
\begin{aligned}
Y&=\beta u+\varepsilon_y,\\
X&=\rho Y+\alpha u+\varepsilon_x.
\end{aligned}\tag{7.A.11}
\end{equation}
定义
\[
e=\frac{X-\rho Y-\alpha u}{\sqrt{1-\rho^2}}.
\]
令 $P_y(u)$ 为随机变量 $Y$ 的密度函数，$P_e(u)$ 为变量 $e$ 的密度函数。定义 $G_y(u)=\log P_y(u)$、$G_e(u)=\log P_e(u)$。类似地，可以得到式（7.A.11）中 $X\leftarrow Y$ 的LiNGAM的对数似然：
\begin{equation}
l(X\leftarrow Y)=\sum_{i=1}^n\left[G_y(y_i)+G_e\left(
 \frac{x_i-\rho y_i-\alpha u_i}{\sqrt{1-\rho^2}}\right)\right]
 -\frac n2\log(1-\rho^2).\tag{7.A.12}
\end{equation}
因此，平均对数似然比为
\begin{equation}
\begin{aligned}
R={}&\frac1n l(X\to Y)-\frac1n l(Y\to X)\\
={}&\frac1n\sum_{i=1}^n\biggl[G_x(x_i)+G_d\left(
 \frac{y_i-\rho x_i-\beta u_i}{\sqrt{1-\rho^2}}\right)\\
&\hspace{4em}{}-G_y(y_i)G_e\left(
 \frac{x_i-\rho y_i-\alpha u_i}{\sqrt{1-\rho^2}}\right)\biggr].
\end{aligned}\tag{7.A.13}
\end{equation}
可以用Taylor展开近似对数密度函数。$G_d\bigl((y_i-\rho x_i-\beta u_i)/\sqrt{1-\rho^2}\bigr)$ 的一阶近似为
\begin{equation}
G_d\left(\frac{y_i-\rho x_i-\beta u_i}{\sqrt{1-\rho^2}}\right)
 =G(y)-(\rho x+\beta u)g(y)+O(\rho^2),\tag{7.A.14}
\end{equation}
其中 $g$ 是 $G$ 的导数。类似地，有
\begin{equation}
G_e\left(\frac{x_i-\rho y_i-\alpha u_i}{\sqrt{1-\rho^2}}\right)
 =G(x)-(\rho y+\alpha u)g(x)+O(\rho^2).\tag{7.A.15}
\end{equation}
将式（7.A.14）、（7.A.15）代入式（7.A.13），得到
\begin{equation}
\begin{aligned}
R\approx{}&\frac1n\sum_{i=1}^n\bigl[G(x_i)+G(y_i)
 -(\rho x_i+\beta u_i)g(y_i)\\
&\hspace{5em}{}-G(y_i)+(\rho y_i+\alpha u_i)g(x_i)\bigr]\\
={}&\frac1n\sum_{i=1}^n\bigl[-(\rho x_i-\beta u_i)g(y_i)
 +(\rho y_i+\alpha u_i)g(x_i)\bigr].
\end{aligned}\tag{7.A.16}
\end{equation}
Hyvärinen and Smith（2013）使用下列逻辑分布密度函数近似对数概率密度函数 $G$：
\begin{equation}
G(z)=-2\log\cosh\left(\frac{\pi}{2\sqrt3}z\right)+\mathrm{const},\tag{7.A.17}
\end{equation}
其中
\[
\cosh(z)=\frac{e^z+e^{-z}}2.
\]
若假设 $\pi/(2\sqrt3)\approx1$，并忽略式（7.A.17）中的常数2，则 $G(z)$ 的导数可近似为
\begin{equation}
g(z)=\frac{dG(z)}{dz}\approx-\frac{e^z-e^{-z}}{e^z+e^{-z}}
 =-\tanh(z),\tag{7.A.18}
\end{equation}
其中 $\tanh(z)=(e^z-e^{-z})/(e^z+e^{-z})$ 为双曲正切函数。

将式（7.A.18）代入式（7.A.16），得到
\begin{equation}
\begin{aligned}
R\approx{}&\frac1n\sum_{i=1}^n\bigl[-(\rho x_i+\beta u_i)g(y_i)
 +(\rho y_i+\alpha u_i)g(x_i)\bigr]\\
={}&\frac1n\sum_{i=1}^n\bigl[(\rho x_i+\beta u_i)\tanh(y_i)
 -(\rho y_i+\alpha u_i)\tanh(x_i)\bigr]\\
={}&\frac1n\sum_{i=1}^n\bigl[\rho\bigl(x_i\tanh(y_i)-y_i\tanh(x_i)\bigr)\\
&\hspace{5em}{}+\beta u_i\tanh(y_i)-\alpha u_i\tanh(x_i)\bigr]\\
={}&\rho\widehat E[x\tanh(y)-y\tanh(x)]
 +\beta\widehat E[u\tanh(y_i)]-\alpha\widehat E[u\tanh(x)].
\end{aligned}\tag{7.A.19}
\end{equation}
定义
\begin{equation}
R_{NC}=\rho\widehat E[x\tanh(y)-y\tanh(x)]\tag{7.A.20}
\end{equation}
以及
\begin{equation}
R_C=\beta\widehat E[u\tanh(y_i)]-\alpha\widehat E[u\tanh(x)].\tag{7.A.21}
\end{equation}
于是，对数似然比可近似为
\begin{equation}
R\approx R_{NC}+R_C.\tag{7.A.22}
\end{equation}
接下来，讨论如何近似四个对数概率密度函数 $G_x,G_y,G_d,G_e$。由中心极限定理，有
\begin{equation}
\frac1n\sum_{i=1}^nG_x(x_i)\to E[G_x(x)]
 =\int P_x(x)\log P_x(x)\,dx=-H(x),\tag{7.A.23}
\end{equation}
其中 $H$ 表示微分熵（differential entropy）。类似地，有
\begin{equation}
\frac1n\sum_{i=1}^nG_d\left(\frac{y_i-\rho x_i-\beta u_i}{\sqrt{1-\rho^2}}\right)
 \to-H\left(\frac{\widehat\varepsilon_y}{\sigma_{\varepsilon_y}}\right),\tag{7.A.24}
\end{equation}
其中
\[
\varepsilon_Y=Y-\rho X-\beta u,\qquad
\sigma_{\varepsilon_Y}=\sqrt{\operatorname{var}(\varepsilon_Y)},
\]
\begin{equation}
\frac1n\sum_{i=1}^nG_y(y_i)\to-H(y),\tag{7.A.25}
\end{equation}
\begin{equation}
\frac1n\sum_{i=1}^nG_e\left(\frac{x_i-\rho y_i-\alpha u_i}{\sqrt{1-\rho^2}}\right)
 \to-H\left(\frac{\widehat\varepsilon_x}{\sigma_e}\right),\tag{7.A.26}
\end{equation}
其中
\[
\varepsilon_x=X-\rho Y-\alpha u,\qquad
\sigma_{\varepsilon_x}=\sqrt{\operatorname{var}(\varepsilon_x)}.
\]
因此，结合式（7.A.13）和式（7.A.23）—（7.A.26），得到
\begin{equation}
R\to-H(x)-H\left(\frac{\widehat\varepsilon_y}{\sigma_{\varepsilon_y}}\right)
 +H(y)+H\left(\frac{\widehat\varepsilon_x}{\sigma_e}\right).\tag{7.A.27}
\end{equation}
定义
\[
\varepsilon_x^0=X-\rho Y,\qquad \varepsilon_y^0=Y-\rho X.
\]
于是，有
\begin{equation}
\varepsilon_x=\varepsilon_x^0-\alpha u\tag{7.A.28}
\end{equation}
以及
\begin{equation}
\varepsilon_y=\varepsilon_y^0-\beta u.\tag{7.A.29}
\end{equation}
由于随机变量 $X,Y$ 的密度函数未知，在实践中计算微分熵并不容易。一种近似微分熵的方法（Hyvärinen 1998；Hyvärinen and Smith 2013）为
\begin{equation}
\widetilde H(z)=H(u)-k_1\{E[\log\cosh(z)]-\gamma\}^2
 -k_2\left\{E\left[ze^{-z^2/2}\right]\right\}^2,\tag{7.A.30}
\end{equation}
其中 $H(u)=\tfrac12(1+2\pi)$、$k_1\approx79.047$、$k_2\approx7.4129$、$\gamma\approx0.37457$。

对 $\widetilde H(z)$ 求导，得到
\begin{equation}
\begin{aligned}
h(z)=\frac{d\widetilde H(z)}{dz}
 ={}&-2k_1\{E[\log\cosh(z)]-\gamma\}E[\tanh(z)]\\
&{}-2k_2E\left[ze^{-z^2/2}\right]
 E\left[(1-z^2)e^{-z^2/2}\right].
\end{aligned}\tag{7.A.31}
\end{equation}
现在利用Taylor展开近似式（7.A.27）中的 $R$。$H\bigl(\widehat\varepsilon_y/\sigma_{\varepsilon_y}\bigr)$ 的一阶Taylor近似为
\begin{equation}
\begin{aligned}
H\left(\frac{\widehat\varepsilon_y}{\sigma_{\varepsilon_y}}\right)
 &=H\left(\frac{\varepsilon_y^0}{\sigma_{\varepsilon_y}}
 -\frac{\beta}{\sigma_{\varepsilon_y}}u\right)\\
 &\approx H\left(\frac{\varepsilon_y^0}{\sigma_{\varepsilon_y}}\right)
 -\frac{\beta}{\sigma_{\varepsilon_y}}
 h\left(\frac{\varepsilon_y^0}{\sigma_{\varepsilon_y}}\right)E[u].
\end{aligned}\tag{7.A.32}
\end{equation}
类似地，有
\begin{equation}
H\left(\frac{\widehat\varepsilon_x}{\sigma_e}\right)
 \approx H\left(\frac{\varepsilon_x^0}{\sigma_{\varepsilon_x}}\right)
 -\frac{\alpha}{\sigma_{\varepsilon_x}}
 h\left(\frac{\varepsilon_x^0}{\sigma_{\varepsilon_x}}\right)E[u].\tag{7.A.33}
\end{equation}
将式（7.A.32）、（7.A.33）代入式（7.A.27），得到
\begin{equation}
\begin{aligned}
R\approx{}&-H(X)-H\left(\frac{\varepsilon_y^0}{\sigma_{\varepsilon_y}}\right)
 +\frac{\beta}{\sigma_{\varepsilon_y}}
 h\left(\frac{\varepsilon_y^0}{\sigma_{\varepsilon_y}}\right)E[u]\\
&{}+H(y)+H\left(\frac{\varepsilon_x^0}{\sigma_{\varepsilon_x}}\right)
 -\frac{\alpha}{\sigma_{\varepsilon_x}}
 h\left(\frac{\varepsilon_x^0}{\sigma_{\varepsilon_x}}\right)E[u]\\
={}&R_0+\left[\frac{\beta}{\sigma_{\varepsilon_y}}
 h\left(\frac{\varepsilon_y^0}{\sigma_{\varepsilon_y}}\right)
 -\frac{\alpha}{\sigma_{\varepsilon_x}}
 h\left(\frac{\varepsilon_x^0}{\sigma_{\varepsilon_x}}\right)\right]E[u],
\end{aligned}\tag{7.A.34}
\end{equation}
其中
\begin{equation}
R_0=-H(x)-H\left(\frac{\varepsilon_y^0}{\sigma_{\varepsilon_y}}\right)
 +H(y)+H\left(\frac{\varepsilon_x^0}{\sigma_e}\right).\tag{7.A.35}
\end{equation}
